\section{A 269-class certificate for the ramified quartic reflection field}
\label{sec:cubic-ramified-quartic-class-certificate}

Put $p=269$, let $t$ be the real cube root of$p$, and put
\[
 z=\zeta_5,\qquad w=\sqrt p\,(z-z^{-1}),\qquad
 F=\mathbb Q(w),\quad L=\mathbb Q(t,w),\quad
 K=\mathbb Q(\sqrt{-3}),\quad M=LK.
 \tag{RQC.1}
\]
Let $\psi$ be the primitive quartic Dirichlet character modulo5 with
$\psi(2)=i$, let $\tau=(\,\cdot\,/269)$, and set $\chi=\psi\tau$.
This is the quartic field different from the cyclotomic one in
AC-G129.  Its additional ramification at269 changes both the number
of roots of unity and the class-formula constant.

\begin{theorem}[AC-G133-quartic-class: the ramified quartic class sector]
\label{AC-G133-quartic-class}
The ideal class group of$L$ has no269-torsion.  Consequently the
nontrivial cubic-character sector of$\operatorname{Cl}(M)[269]$ is zero.
The proof does not assert a full class number, a complete unit basis,
vanishing of the trivial cubic sector of$\operatorname{Cl}(M)[269]$,
or any subsequent deformation-host or model statement.
\end{theorem}

\begin{proof}
\proofstep{1. The precise quartic field and its local data.}
The fields $\mathbb Q(z)$ and $\mathbb Q(\sqrt p)$ are disjoint:
the unique quadratic subfield of the former is$\mathbb Q(\sqrt5)$.
In their compositum let$s$ send$z$ to$z^2$ and fix$\sqrt p$, and
let$r$ fix$z$ and negate$\sqrt p$.  Then$s^2w=rw=-w$, whereas
$sw\ne\pm w$.  Thus the stabilizer of$w$ is$\{1,s^2r\}$.
This is exactly the kernel of$\psi\tau$ on$C_4\times C_2$.
It follows that$F/\mathbb Q$ is cyclic of degree four, with character
group$\{1,\psi^2,\chi,\bar\chi\}$, and
\[
 w^4+5p w^2+5p^2=0,\qquad
 w^2=-p(5\pm\sqrt5)/2.
 \tag{RQC.2}
\]
The character$\chi$ is primitive of conductor$5p$ and is odd;
its square is the quadratic character$\delta=\psi^2$ of conductor5.
In particular$F$ is totally imaginary and unramified at3.

The polynomial$X^3-p$ is Eisenstein at$p$.  At3, write a cube root
of$p$ in the form$-(1+3v)$: the required equation is
$v+3v^2+3v^3=-30$, whose derivative is1 modulo3.
Hensel's lemma therefore gives a root in$\mathbb Q_3$, and
\[
 \mathbb Q(t)\otimes\mathbb Q_3
      =\mathbb Q_3\times\mathbb Q_3(\sqrt{-3}).
\]
The cubic discriminant exponents are one at3 and two at$p$;
there is no other ramification by the polynomial discriminant.
The degrees three and four make$\mathbb Q(t)$ and$F$ disjoint.
The local quartic ramification indices are four at5 and two at$p$,
as follows directly from the two primitive local factors of$\chi$.
Both are tame.  At$p$ its ramified quadratic extension and the
tamely ramified cubic extension have compositum ramification index6.
The total degree is twelve, so the tame different contributes
$12(1-1/6)=10$ to the discriminant exponent there.  At5 the cubic
algebra is unramified and contributes three copies of the quartic
different, for exponent9.  At3 the unramified quartic base change
multiplies the cubic discriminant exponent by four.  Hence
\[
 (r_1(L),r_2(L))=(0,6),\qquad
 D_L=3^4 5^9p^{10}
       =313861364090558493911019298828125.
 \tag{RQC.3}
\]
This computation does not use a polynomial-order discriminant as
though it were the field discriminant.

The splitting field of$X^3-p$ has group$S_3$ and quadratic subfield$K$.
Its intersection with the abelian field$F$ is$\mathbb Q$: otherwise
it would contain$K$, which is ramified at3.  Thus
$\operatorname{Gal}(M/\mathbb Q)=S_3\times C_4$, and$L$ is the fixed
field of one reflection in$S_3$.  That reflection normally generates
$S_3$, so every abelian subextension of$L$ lies in$F$.
Any root of unity in$L$ therefore lies in$F$.
The only quadratic subfield of$F$ is$\mathbb Q(\sqrt5)$; the possible
cyclotomic degrees dividing four exclude all roots of unity other than
$\pm1$, except potentially fifth roots.  Those would make
$F=\mathbb Q(\zeta_5)$, contrary to its ramification at$p$.
Consequently
\[
                         |\mu_L|=2.
 \tag{RQC.4}
\]

\proofstep{2. Exact units and their269-saturation.}
Use the ordered basis$t^a w^b$, with$0\le a<3$, $0\le b<4$,
in lexicographic order.  The five units below have denominators
\[
                    (d_1,d_2,d_3,d_4,d_5)=(269,1,807,434166,434166).
 \tag{RQC.5}
\]
For each$i$, read the following coefficient array row by row to
obtain its twelve numerators, and divide by$d_i$.
\begingroup\scriptsize
\[
 U_1=\begin{pmatrix}
538&0&1\\
0&0&0\\
0&0&0\\
0&0&0
\end{pmatrix}.
\]
\[
 U_2=\begin{pmatrix}
270&0&0\\
0&-739&0\\
0&0&108\\
0&0&0
\end{pmatrix}.
\]
\[
 U_3=\begin{pmatrix}
56040048004&0&150729577\\
0&8680396508&0\\
23348834&0&1344189772\\
0&3616480&0
\end{pmatrix}.
\]
\[
 U_4=\begin{pmatrix}
-55471141244894929837391250836&-7114308155340666243833892654&-56995709003186289184103540\\
-7309838007814710982052088&-14105692072688439819953489911&937014432047789939511734799\\
-14493373250324747359323571&962767371929317496360784&3516294549348002120183821576\\
25571326110405958286987679&3612936414800857297005739&26274127950187693604979
\end{pmatrix}.
\]
\[
 U_5=\begin{pmatrix}
-21188087364390629115228010464&-1679460225000426189781561242&-56995709003186289184103540\\
-4517728189718195097835902&-5387894948998345378336713084&221199065653505753324695437\\
-14493373250324747359323571&595022910412010111333019&1343104928559150944288897379\\
6036574145212979063569527&3612936414800857297005739&16238308009682117277954
\end{pmatrix}.
\]
\endgroup
These rational arrays define$u_i$ exactly.  The multiplication laws
are$t^3=p$ and$w^4=-5pw^2-5p^2$.
For each$u_i$ the certificate forms the twelve rational multiplication
columns and computes the monic characteristic polynomial.  Every
coefficient is integral and the constant coefficient is1; direct
Horner evaluation in the same rational ring gives zero.
Thus$u_i$ and$u_i^{-1}$ are integral, so they are global units.
All five complete polynomials and their exact rational identities
are preserved in\texttt{all\_units\_exact\_certificate.json}, SHA256
\texttt{16d44f9af0d17bac22dd703aaf8d75a4ae8}\allowbreak
\texttt{1f2496dc8b6e31cf008207f5ba878}.
The separate implementation\texttt{verify\_finite\_independently.py}
recomputes all five annihilations by reducing a bivariate product,
without using the candidate program's multiplication matrices.
The raw candidate coordinate was$w-t$; the reduced polynomial
$X^{12}+13450X^6+9045125$ has the verified generator$-w/t$.
Neither that reduced coordinate nor an approximate embedding is
substituted for the exact arrays above.

For saturation take$\ell=40351$, primitive root3, and
$h=3^{150}=10281\pmod\ell$, which has order269.
Complete trial division checks primality; the powers for the prime
divisors$2,3,5,269$ of$\ell-1$ check primitivity.
The first five reduction maps are given by
\[
 (t,w)=(21085,12062),(21085,14774),(21085,25577),
       (21085,28289),(23039,12062).
 \tag{RQC.6}
\]
They satisfy both defining equations with nonzero derivatives modulo
$\ell$.  Hensel lifts give embeddings$L\hookrightarrow\mathbb Q_\ell$;
every global algebraic integer maps into$\mathbb Z_\ell$.
They are therefore genuine maps of$\mathcal O_L$, even without an
assumption that the displayed polynomial order is maximal.
For the image$\bar u_i$ under map$j$, write
$\bar u_i^{150}=h^{c_{ji}}$.  Exact powering gives
\[
 (c_{ji})=
 \begin{pmatrix}
 190&49&187&27&174\\
 79&49&82&174&55\\
 79&49&82&13&27\\
 190&49&187&55&13\\
 190&142&148&17&233
 \end{pmatrix},\qquad
 \det(c_{ji})=-7668928260\equiv3\pmod{269}.
 \tag{RQC.7}
\]
The certificate checks all twelve residue maps, not only this minor.
Each resulting additive logarithm kills269th powers and$\mu_L$.
The nonsingular minor proves independence modulo269.
Any integral relation modulo torsion would have all its exponents
divisible by269; division and repetition contradict a nonzero
relation.  The units are therefore multiplicatively independent.
The unit rank is five, so
\[
 U=\mu_L\langle u_1,\ldots,u_5\rangle,\qquad
 m=[\mathcal O_L^\times:U]<\infty,\qquad269\nmid m.
 \tag{RQC.8}
\]
For the last assertion, an element of order269 in the quotient would
give$v^{269}=\xi\prod u_i^{a_i}$.  Reduction makes every$a_i$ divisible
by269; after dividing these powers, $v$ differs from an element of$U$
by a root of unity, since its269th power is torsion.  This is a
contradiction.  No completeness of the candidate unit basis is used.

\proofstep{3. A finite bound for the unknown unit index.}
For a degree$n$ number field put
\[
 P_n=\prod_{j=1}^{n}\bigl(4\tbinom nj+1\bigr),\qquad
 \delta_n=\frac{\log2}{n(P_n+1)}.
 \tag{RQC.9}
\]
If every conjugate of a nontorsion unit had modulus at most
$e^{\delta_n}$, the characteristic polynomials of its first
$P_n+1$ powers would have integral$j$th coefficient bounded by
$2\binom nj$.  There are at most$P_n$ such polynomials.  Equality
of two of them makes the largest conjugate modulus equal to1.
The unit norm then makes all moduli1.  All positive powers would
again have bounded integral polynomials, giving only finitely many
possible roots; two powers would coincide, a contradiction.
Thus some conjugate has log modulus greater than$\delta_n$.

Project the weighted logarithm lattice to its first$r=r_1+r_2-1$
coordinates.  Its last coordinate is minus their sum.
No nonzero lattice vector can have all retained coordinates of
absolute value less than$\delta_n/r$: otherwise every unweighted
log modulus would have absolute value less than$\delta_n$.
The open cube of side$\delta_n/r$ injects almost everywhere into
the quotient torus, since a difference lies in the forbidden cube.
Comparing volumes gives$R_L\ge(\delta_n/r)^r$.
For$n=12$, $r=5$ and$\log2>1/2$ this yields
\[
 \begin{gathered}
 P=P_{12}=95338564630654861425388299414125,\\
 R_L>[120(P+1)]^{-5}.
 \end{gathered}
 \tag{RQC.10}
\]
Representatives for the six complex places are
\[
 t\longmapsto p^{1/3}\omega^j\ (j=0,1,2),\qquad
 w\longmapsto i\sqrt{p(5\pm\sqrt5)/2}.
 \tag{RQC.11}
\]
The regulator$R_U$ is the absolute determinant of any five rows
of$2\log|\sigma(u_i)|$.  The900-bit ball evaluation of these
exact expressions gives
\[
 150862773<R_U<150862774,\qquad
 m=R_U/R_L<M_*:=150862774[120(P+1)]^5.
 \tag{RQC.12}
\]
All six rows are retained in the output, with each column sum
containing zero.  These are the projected weighted logarithms,
not an unnormalized Euclidean hyperplane covolume.

\proofstep{4. The exact Artin factors and Dirichlet constants.}
Let$\rho$ be the standard representation of the$S_3$ splitting field
of$X^3-p$.  The permutation representation for$L$ is
$(1\oplus\rho)\otimes(1\oplus\delta\oplus\chi\oplus\bar\chi)$.
Artin formalism, including the ramified inertia invariants, gives
\[
 \operatorname*{res}_{s=1}\zeta_L(s)
  =L(1,\delta)|L(1,\chi)|^2
        L(1,\rho)L(1,\rho\delta)|L(1,\rho\chi)|^2.
 \tag{RQC.13}
\]
Write$\varphi=(1+\sqrt5)/2$.
The four-term finite Fourier identity modulo5, applied to
$-\log|1-re^{i\theta}|=\sum_{n\ge1}r^n\cos(n\theta)/n$, gives
\[
 L(1,\delta)
   =-\frac1{\sqrt5}\sum_{a=1}^4\delta(a)
                   \log(2\sin(\pi a/5))
   =\frac{2\log\varphi}{\sqrt5}.
 \tag{RQC.14}
\]
The limit$r\uparrow1$ uses Dirichlet convergence and Abel's theorem;
the last equality uses$\sin(2\pi/5)/\sin(\pi/5)=\varphi$.

The exact finite Bernoulli sum for the primitive odd character
of conductor$q=1345$ is
\[
 B_{1,\chi}=\frac1{1345}\sum_{a=1}^{1344}a\psi(a)\tau(a)=6+10i.
 \tag{RQC.15}
\]
This is checked with integer arithmetic: the real and imaginary
numerator sums are respectively8070 and13450.  The quadratic table
is computed both by Euler's criterion and independently as the set
of nonzero squares modulo269.  No approximate character values occur.
The primitive odd Dirichlet functional equation and the Gauss norm
$|\mathfrak g(\chi)|=\sqrt q$ give
\[
 |L(1,\chi)|^2=\frac{\pi^2}{1345}|B_{1,\chi}|^2
                   =\frac{136\pi^2}{1345}.
 \tag{RQC.16}
\]
The phase is unnecessary for this absolute square.  Consequently the
first two factors in (RQC.13) have the exact product
$\bigl(2\log\varphi/\sqrt5\bigr)\bigl(136\pi^2/1345\bigr)$.

Let$L(s,\rho)=\sum a_n n^{-s}$.  Its multiplicative coefficients
have, for a prime$\ell\ne3,p$, initial value
$a_\ell=\#\{x\bmod\ell:x^3=p\}-1$ and recurrence
\[
 a_{\ell^{v+1}}=a_\ell a_{\ell^v}
                   -(-3/\ell)a_{\ell^{v-1}},\qquad
 a_{\ell^0}=1,\quad a_{3^v}=1,\quad a_{p^v}=0\quad(v\ge1).
 \tag{RQC.17}
\]
For$\ell\equiv2\pmod3$ the cubic map is bijective, and the local
coefficients are1 in even degree and0 in odd degree.
For$\ell\equiv1\pmod3$, testing$p^{(\ell-1)/3}\bmod\ell$ gives
the three-root or zero-root case.  At3 the local cubic extension
is trivial, as proved above; at$p$ cubic inertia has no standard
invariants.  Thus the two exceptional factors in (RQC.17) are exact.
The twisted coefficients are$a_n\delta(n)$ and
$a_n\psi(n)\tau(n)$, with their ordinary zero extensions at ramified
primes.  Their primitive conductors are
\[
 Q_0=3p^2=217083,\qquad Q_5=75p^2=5427075.
 \tag{RQC.18}
\]
Indeed$\rho=\operatorname{Ind}_{K}^{\mathbb Q}\nu$ for the cubic
Kummer character of conductor$(p)$.  At$p$, multiplying its order-three
inertia character by the nontrivial quadratic norm pullback of$\tau$
gives an order-six character, still of conductor exponent one over$K$.
At5 the nontrivial norm pullback of$\psi$ or$\delta$ has conductor$(5)$.
The norm maps on the residue fields are surjective; the primes5 and$p$
are both inert in$K$.  No cancellation occurs.  At3 all these norm
characters are unramified.  Induction multiplies the conductor norm
by$|D_K|=3$, proving (RQC.18), including the shared ramification at$p$.

\proofstep{5. The complete theta normalization and certified phase.}
The ring$\mathcal O_K=\mathbb Z[\omega]$ has class number one
(the imaginary quadratic Minkowski bound is$2\sqrt3/\pi<2$,
so an integral ideal representative in each class has norm one).
Each character needed in (RQC.18) is a primitive finite character
$\nu_f$ modulo$f\mathcal O_K$, with$f=p$ or$5p$, trivial on$\mu_6$.
For the cubic character, triviality on the cubic roots of unity follows
from$9\mid p+1$; its order also kills$-1$.  Norm pullbacks are
trivial on all six units automatically.
Extend$\nu_f$ by zero on nonunits and put
\[
 a=\frac{2\pi}{\sqrt3 f}=\frac{2\pi}{\sqrt Q},\qquad
 \Theta(t)=\frac16\sum_{\alpha\in\mathcal O_K}
             \nu_f(\alpha)e^{-a|\alpha|^2t}
           =\sum_{n\ge1}b_n e^{-ant}.
 \tag{RQC.19}
\]
The second identity counts the six generators of each principal ideal.
The finite additive Fourier transform of$\nu_f$ at a unit frequency
is a fixed Gauss constant times its inverse character, by substitution.
It vanishes at a nonunit frequency: choose a unit in the corresponding
residue-field factor on which the multiplicative character is nontrivial
but which fixes the additive character, and use the identity in the
other factors.  Primitivity supplies this unit at every factor of the
squarefree ideal$(f)$.  Additive orthogonality then gives Gauss modulus$f$.

For completeness the Euclidean lattice$f\mathcal O_K$ has area
$f^2\sqrt3/2$, and its dual is$2/(f\sqrt{-3})\mathcal O_K$.
The Fourier transform of$e^{-at|x|^2}$ is
$\pi/(at)\,e^{-\pi^2|\xi|^2/(at)}$.
Write$d=2/\sqrt{-3}$ and a dual vector$d\beta/f$.
The residue Fourier sum is$G\bar\nu_f(\bar\beta)$, with its fixed
unit factor absorbed in$G$.  Conjugating$\beta$ preserves its norm.
The lattice constants give$1/(ft)$ and change the Gaussian exponent
to$a|\beta|^2/t$.  Poisson summation therefore gives exactly
\[
       \Theta(t)=\varepsilon t^{-1}\bar\Theta(1/t),\qquad
       \varepsilon=G/f,\qquad |\varepsilon|=1.
 \tag{RQC.20}
\]
Both sides have no constant term; the nonprincipal character also
makes the transform decay exponentially at zero.
Splitting the Mellin integral at1 and integrating each exponentially
convergent term gives
\[
 L(1)=\sum_{n\ge1}\frac{b_n e^{-an}}n
       +\varepsilon a\sum_{n\ge1}\bar b_n E_1(an),\qquad
 E_1(x)=\int_1^\infty e^{-xt}\frac{dt}{t}.
 \tag{RQC.21}
\]
The interval evaluation proves$\Theta(1)\ne0$ for each of the three
characters.  Thus the actual phase is determined, without a convention
guess, by
\[
                     \varepsilon=\Theta(1)/\overline{\Theta(1)}.
 \tag{RQC.22}
\]
For the two real coefficient sequences it is1.  The complex quotient
in (RQC.22) is evaluated as a ball for the quartic twist.

\proofstep{6. Absolute tails and rational isolation.}
Every local coefficient in (RQC.17) has modulus at most$v+1$;
hence$|b_n|\le d(n)\le n$ for every twist.
For$z_a=e^{-a}$ and$H=e^{-a(N+1)}$, the omitted absolute tails
of the first sum in (RQC.21), the second sum before multiplication
by$a$, and the theta series are respectively bounded by
\[
 \frac H{1-z_a},\qquad
 \frac H{a(1-z_a)},\qquad
 \frac{H((N+1)-Nz_a)}{(1-z_a)^2}.
 \tag{RQC.23}
\]
The second bound uses$E_1(x)\le e^{-x}/x$.
The code sums by residues modulo5, with and without the additional
factor$\tau(n)$.  Each residue subtotal gets the full positive tail
in (RQC.23); this only enlarges the safe error radius.
The same exponential and$E_1$ values are reused for these two
sequences, without cancelling any of their certified errors.
The cutoffs are$N_0=40000$ and$N_5=200000$, at900-bit precision.
All operations and additions of tail radii use inclusion arithmetic.

The analytic class-number formula with (RQC.3)--(RQC.4) is
\[
 h_LR_L=\frac{2\sqrt{D_L}}{(2\pi)^6}
                \operatorname*{res}_{s=1}\zeta_L(s).
 \tag{RQC.24}
\]
Combining (RQC.13)--(RQC.24) with the ball determinant in (RQC.12)
gives the strict rational-isolation comparison
\[
             \left|\frac{h_LR_L}{R_U}-1428\right|<\frac1{M_*}.
 \tag{RQC.25}
\]
In particular the strict ball bounds are
\[
 \left|h_LR_L/R_U-1428\right|<7\cdot10^{-202},\qquad
                   1/M_*>3.38\cdot10^{-179}.
\]
The exact successful interval and all inputs are preserved in
\texttt{d269\_analytic\_interval.json}; the verifier is
\texttt{certify\_d269\_analytic.py}.  The comparison in (RQC.25)
is made on ball objects with outward-rounded tails, not printed digits.
The integer1428 is only the preselected target for this inequality.

Since$R_U=mR_L$, the quantity being enclosed is$h_L/m$.
If$h_L\ne1428m$, the nonzero integer$h_L-1428m$ would give
$|h_L/m-1428|\ge1/m>1/M_*$, contradicting (RQC.25).
Thus$h_L=1428m$.  Both1428 and$m$ are prime to269, by (RQC.8),
so$\operatorname{Cl}(L)[269]=0$.
The equality$m=1$ is neither needed nor asserted.

Finally let$s$ be the reflection fixing$L$ in$M$.
Extension and norm on ideal classes have composites2 and$1+s$.
Because2 is invertible on the269-primary part, its$s$-invariant
subspace is zero.  Hence$s$ acts as$-1$, as do all its conjugates.
Products of conjugate reflections are the cubic rotations and act
as1.  The nontrivial cubic-character sector is therefore zero.
This does not annihilate the trivial cubic sector and proves only
the stated class conclusions.
\end{proof}

\paragraph{Source and certification scope.}
The analytic class-number formula and projected weighted-log convention
are the statement on printed p.1 and definition on printed p.4 of
A.~Sutherland, MIT18.785 Lecture19 (13 November2019),
\url{https://math.mit.edu/classes/18.785/2019fa/LectureNotes19.pdf}.
Those passages and the surrounding measure convention were freshly
read in the immediately preceding G129/G131 reviews and applied here;
no new complete proof of the classical theorem is claimed.
The primitive Dirichlet functional equation and special-value formula
are DLMF25.15.3,25.15.5--6 at\url{https://dlmf.nist.gov/25.15}.
The specialized finite Fourier, Poisson, phase, Mellin, regulator bound
and rational-isolation arguments are proved explicitly above.
The inclusion-arithmetic API is python-flint0.9.0, including
\texttt{arb.expint} at order1, with its documented outward radii:
\url{https://python-flint.readthedocs.io/en/latest/arb.html}.
No full runtime implementation proof, Lean execution, GRH assumption,
or provisional PARI class-group completeness is used.

\paragraph{Geometric dictionary and exact remaining scope.}
The conclusion controls the indicated cubic sector of unramified
abelian covers of$\operatorname{Spec}\mathcal O_M$, via the retained
class-field and finite-etale-cover dictionary (SGA1, ExposeV,
sections4 and7).  Reflection descent is exact prime-to269 norm and
extension.  No new purity statement or automorphic comparison is
deduced.  Combining this finite certificate with a different field,
a changed residual tuple, or a native deformation host requires the
corresponding separate arithmetic and source interfaces.
